\scanpage{114}
\mspage{NA 27}
% archive sheet number: 113
% Typescript upside down in the scan; transcribed in its reading direction.

Page 1074

($^{**}$) (8 mai 1986) En écrivant ici ``du\deleted{\unsure{un}} cru de Kashiwara'', j'avais à l'esprit, bien sûr, l'incident au séminaire Goulaouic-Schwartz, telle que me l'avait rapporté Mebkhout. Je n'ai plus lieu de croire à présent qu'en cette occasion, Kashiwara ait agi de façon incorrecte, \deleted{\unsure{et pensé|et pense} que Mebkhout} Par contre, j'estime que d'autres actes d'appropriation du cru de Kashiwara et de certains de ses collaborateurs, dont il a été question par ailleurs, méritent bel et bien le nom d'``escroqueries''. J'avais pris soin d'ailleurs de faire parvenir à Kashiwara \deleted{dès} un exemplaire de ReS IV dès après tirage, suivi le 9 février dernier par une lettre d'excuses au sujet des déformations de la version Mebkhout, dont je m'étais fait l'écho\inserted{\handwritten{sans réserves}}. J'ai appris par ailleurs qu'il a reçu l'un et l'autre envoi. \deleted{\unsure{Il n'|Il m'}}\handwritten{\unsure{,}} Je n'ai reçu à ce jour aucun\deleted{\unsure{s}} signe de vie de lui, ni d'explication au sujet de\deleted{s} certains faits exposés dans Récoltes et Semailles le concernant, et qui ne peuvent guère faire l'objet \deleted{\unsure{d'un}}de doutes.

Page 1076

($^{*}$) (8 mai 1986) J'exprime ici vis-à-vis de Kashiwara un préjugé défavorable (à dire le moins) qui s'est avéré injustifié: la \unsure{vertité|vérité}, c'est que jusqu'en 1979 ou 1980, p e r s o n n e (me semble-t-il) \deleted{n'y} ne comprenait la place et la portée de ce résultat -- et Kashiwara, qui \deleted{\unsure{le|l'}}avait été le premier à le pressentir, ne le comprenait sans doute pas moins que les autres, y compris même peut-être Mebkhout lui-même. Après ``Verdier en personne'', j'ai supprimé ici la fin de la phrase\inserted{\handwritten{du texte original}}, encore plus désobligeante à l'encontre de Kashiwara que le début, et toute aussi abusive.

Page 1088

($^{*}$) (8 mai 1986) Je rappelle que c'est là la version de Mebkhout, qui à présent m'apparaît comme affabulatoire. Du coup, cela modifie de façon draconienne les présupposés ( m e s présupposés) dans l'incident au séminaire Goulaouic-Schwartz, dont il va être question. Les réflexions que cet incident m'avait inspirées l'an dernier ne m'en paraissent pourtant pas moins pertinentes. Et il n'y a aucun doute pour moi que l'avachissement et l'abdication de tous (ou du très grand nombre), devant la dégradation de l'éthique scientifique, est pour beaucoup dans l'exacerbation de dispositions conflictuelles à outrance, comme celles qui ont opposé Mebkhout à Kashiwara\deleted{\unsure{,}} et à d'autres.\deleted{\unsure{qu'il compte parmi les amis de celui}\illegible{1}}
